%%% XeLaTeX-article %%%
%# -*- coding: utf-8 -*-
%!TEX encoding = UTF-8 Unicode
%!TEX TS-program = xelatex
%---------------------虽然加了%还是要保留！

\documentclass[17pt]{beamer}
\mode<presentation>
{
\usetheme[width=40pt]{Hannover}
\usecolortheme[]{dove}
\usefonttheme[]{structurebold}
\setbeameroption{hide notes}
}

\usepackage{fontspec}
\usepackage{polyglossia}
\setmainfont{Arial} %设置主字体
\newfontfamily\sanskritfont[Script=Devanagari,Mapping=romantodevanagari,Scale=1.15]{Sanskrit 2003}             %输出天城体
%\newfontfamily\sanskritfont[Mapping=tex-text]{Times New Roman}              %输出转写
\doublehyphendemerits=-10000
\newcommand{\skt}[1]{{\sanskritfont{#1}}} %输出天城体
\newcommand{\skttrans}[1]{{\skt{#1}~#1}}  %输出天城体和转写
%----------------------------------------------------设置梵文输入方法 danda । ॥

\usepackage[UTF8,fontset=windows]{ctex}
\usepackage{amsmath}
%----------------------------------------------------设置中文环境

\usepackage{graphicx}
\usepackage{flafter}
\graphicspath{{pic/}}
\usepackage{booktabs}
\usepackage{nicematrix}
\newenvironment{indentlist}
  {\begin{list}{}{\setlength{\leftmargin}{2em}\setlength{\itemsep}{0.5em}}}
  {\end{list}}
%-----------------------------------------插图表格

\usepackage{hyperref}
\usepackage[dvipsnames]{xcolor}
\usepackage{colortbl}
\definecolor{light-gray}{gray}{0.9}
%------------------------------颜色

\newcommand{\verbroot}[1]{\textcolor{red}{$\sqrt{}$#1}}
\newcommand{\sktroot}[1]{{\verbroot{\skt{#1}}}}
\newcommand{\skttransroot}[1]{{\sktroot{#1}~\textcolor{red}{#1}}}

\newcommand{\nounstem}[1]{\textcolor{red}{#1\nobreakdash-}}
\newcommand{\sktnounstem}[1]{{\textcolor{red}{\skt{#1\nobreakdash-}}}}
\newcommand{\skttransnounstem}[1]{{\sktnounstem{#1}~\nounstem{#1}}}

\newcommand{\verbstem}[1]{\textcolor{blue}{#1\nobreakdash-}}
\newcommand{\sktverbstem}[1]{{\textcolor{blue}{\skt{#1\nobreakdash-}}}}
\newcommand{\skttransverbstem}[1]{{\sktverbstem{#1}~\verbstem{#1}}}

\newcommand{\wordending}[1]{\textcolor{Orange}{\nobreakdash-#1}}
\newcommand{\sktending}[1]{{\textcolor{Orange}{\skt{-#1}}}}
\newcommand{\skttransending}[1]{{\sktending{#1}~\wordending{#1}}}

\newcommand{\fullpada}[1]{\textcolor{OliveGreen}{#1}}
\newcommand{\sktpada}[1]{{\textcolor{OliveGreen}{\skt{#1}}}}
\newcommand{\skttranspada}[1]{{\sktpada{#1}~\fullpada{#1}}}

\newcommand{\pratyaya}[1]{\mbox{\textcolor{Plum}{#1}}}
\newcommand{\sktpratyaya}[1]{{\textcolor{Plum}{\skt{#1}}}}
\newcommand{\skttranspratyaya}[1]{{\sktpratyaya{#1}~\pratyaya{#1}}}

\newcommand{\reconstruction}[1]{\textcolor{gray}{*#1}}
\newcommand{\fullsentence}[1]{\textcolor{MidnightBlue}{#1}}
\newcommand{\sktsentence}[1]{\textcolor{MidnightBlue}{\skt{#1}}}

\newcommand{\veryimportant}[1]{\textcolor{red}{#1}}
\newcommand{\important}[1]{\textcolor{blue}{#1}}
\newcommand{\notsoimportant}[1]{\textcolor{gray}{#1}}
\newcommand{\dicturl}[2]{\href{#1}{\mbox{\texttt{#2}}}}
%-------------------------------------------词根等标颜色

\title{{梵语提高}}
\subtitle{37. 数词}
\author[张雪杉]{文学院~~张雪杉\\zhangxueshan@sdnu.edu.cn}
\date{}

\begin{document}

\begin{frame}
  \titlepage
\end{frame}

\begin{frame}
  \frametitle{本节内容}
  \small
  \tableofcontents
\end{frame}

\section{上节作业}

\begin{frame}{《薄伽梵歌》1.26-29}
  \footnotesize{
    \begin{verse}
      \mbox{\skt{tatrāpaśyatsthitānpārthaḥ pitṝnatha pitāmahān ।}}\\
      \mbox{\skt{ācāryānmātulānbhrātṝnputrānpautrānsakīṃstathā ।। 1-26 ।।}}\\
      \skt{śvaśurānsuhṛdaścaiva senayorubhayorapi ।}\\
      \mbox{\skt{tānsamīkṣya sa kaunteyaḥ sarvānbandhūnavasthitān ।। 1-27 ।।}}\\
      \skt{kṛpayā parayāviṣṭo viṣīdannidamabravīt ।}\\
      \skt{arjuna uvāca}\\
      \mbox{\skt{dṛṣṭvemaṃ svajanaṃ kṛṣṇa yuyutsuṃ samupasthitam ।। 1-28 ।।}}\\
      \skt{sīdanti mama gātrāṇi mukhaṃ ca pariśuṣyati ।}\\
      \mbox{\skt{vepathuśca śarīre me romaharṣaśca jāyate ।। 1-29 ।।}}\\
    \end{verse}
  }
\end{frame}

\begin{frame}{《薄伽梵歌》1.30-33}
  \footnotesize{
    \begin{verse}
      \skt{gāṇḍīvaṃ sraṃsate hastāttvakcaiva paridahyate ।}\\
      \mbox{\skt{na ca śaknomyavasthātuṃ bhramatīva ca me manaḥ ।। 1-30 ।।}}\\
      \skt{nimittāni ca paśyāmi viparītāni keśava ।}\\
      \mbox{\skt{na ca śreyo 'nupaśyāmi hatvā svajanamāhave ।। 1-31 ।।}}\\
      \skt{na kāṅkṣe vijayaṃ kṛṣṇa na ca rājyaṃ sukhāni ca ।}\\
      \mbox{\skt{kiṃ no rājyena govinda kiṃ bhogairjīvitena vā ।। 1-32 ।।}}\\
      \skt{yeṣāmarthe kāṅkṣitaṃ no rājyaṃ bhogāḥ sukhāni ca ।}\\
      \mbox{\skt{ta ime 'vasthitā yuddhe prāṇāṃstyaktvā dhanāni ca ।। 1-33 ।।}}\\
    \end{verse}
  }
\end{frame}

\section{基数词}
\begin{frame}{\insertsection}
  \small
  \tableofcontents[currentsection]
\end{frame}

\subsection{概述}
\begin{frame}{基数词 \insertsubsection}
  \small
  \begin{itemize}
    \item \nounstem{eka} (1) \nounstem{dvi} (2) 
    \nounstem{tri} (3) \nounstem{catur} (4) \\
    是形容词，要按修饰的名词分性变格。
    \item \nounstem{pañca} (5) \nounstem{ṣaṣ} (6) \nounstem{sapta} (7) \\
    \nounstem{aṣṭa} (8) \nounstem{nava} (9) \nounstem{daśa} (10)\\
    按中性名词变格。
  \end{itemize}
\end{frame}

\begin{frame}{基数词 \insertsubsection}
  \small
  \begin{itemize}
    \item \nounstem{eka} (1)：变格类似tad （第22章）
    \\ 三个数都有，双数和复数表示一些、某些
    \item \nounstem{dvi} (2)：只有双数。
    \item \nounstem{tri} (3) \textasciitilde{} \nounstem{daśa} (10) 只有复数。

    \item 更大的查字典\\百：\nounstem{śata} (n.)；千：\nounstem{sahasra} (n.)\\
  \end{itemize}
\end{frame}

\subsection{二}
\begin{frame}{dvi- (2)}
  \small
  \centering
  \begin{NiceTabular}{|c|c|c|c|}[hvlines, rules/width=0.3pt, rules/color=gray]
    & \important{阳性} & \important{中性} & \important{阴性} \\
    主/呼/业格 & \fullpada{dvau} & \fullpada{dve} & \fullpada{dve} \\
    具/为/从格 & \fullpada{dvābhyām} & \fullpada{dvābhyām} & \fullpada{dvābhyām} \\
    属/依格   & \fullpada{dvayoḥ} & \fullpada{dvayoḥ} & \fullpada{dvayoḥ} \\
  \end{NiceTabular}
  \begin{itemize}
    \item 词干 \nounstem{dvi} 是复合词中使用的形式。\\
    \nounstem{dvipad} 两脚的 \nounstem{trilocana} 三眼的
  \end{itemize}
\end{frame}

\subsection{三和四}
\begin{frame}{\mbox{tri- (3) 和 catur- (4) }}
  \footnotesize
  \centering
  \resizebox{\textwidth}{!}{%
  \begin{NiceTabular}{|c|c|c|c|}[hvlines, rules/width=0.3pt, rules/color=gray]
    \important{tri- (3)} & 阳 & 中 & 阴 \\
    主/呼 & \fullpada{trayaḥ} & \fullpada{trīṇi} & \fullpada{tisraḥ} \\
    业格 & \fullpada{trīn} & \fullpada{trīṇi} & \fullpada{tisraḥ} \\
    具格 & \fullpada{tribhiḥ} & \fullpada{tribhiḥ} & \fullpada{tisṛbhiḥ} \\
    为/从格 & \fullpada{tribhyaḥ} & \fullpada{tribhyaḥ} & \fullpada{tisṛbhyaḥ} \\
    属格 & \fullpada{trayāṇām} & \fullpada{trayāṇām} & \fullpada{tisṛṇām} \\
    依格 & \fullpada{triṣu} & \fullpada{triṣu} & \fullpada{tisṛṣu} \\

    \important{catur- (4)} & 阳 & 中 & 阴 \\
    主/呼 & \fullpada{catvāraḥ} & \fullpada{catvāri} & \fullpada{catasraḥ} \\
    业格 & \fullpada{caturaḥ} & \fullpada{catvāri} & \fullpada{catasraḥ} \\
    具格 & \fullpada{caturbhiḥ} & \fullpada{caturbhiḥ} & \fullpada{catasṛbhiḥ} \\
    为/从格 & \fullpada{caturbhyaḥ} & \fullpada{caturbhyaḥ} & \fullpada{catasṛbhyaḥ} \\
    属格 & \fullpada{caturṇām} & \fullpada{caturṇām} & \fullpada{catasṛṇām} \\
    依格 & \fullpada{caturṣu} & \fullpada{caturṣu} & \fullpada{catasṛṣu} \\
  \end{NiceTabular}%
  }
\end{frame}

\subsection{五到十}
\begin{frame}{\mbox{pañca- (5) \textasciitilde{} daśa- (10) }}
  \footnotesize
  \centering
  \resizebox{\textwidth}{!}{%
  \begin{NiceTabular}{|c|c|c|c|}[hvlines, rules/width=0.3pt, rules/color=gray]
     & \important{pañca (5)} & \important{ṣaṣ (6)} & \important{sapta (7)} \\
    主/呼/业 & \fullpada{pañca} & \fullpada{ṣaṭ} & \fullpada{sapta} \\
    具格 & \fullpada{pañcabhiḥ} & \fullpada{ṣaḍbhiḥ} & \fullpada{saptabhiḥ} \\
    为/从格 & \fullpada{pañcabhyaḥ} & \fullpada{ṣaḍbhyaḥ} & \fullpada{saptabhyaḥ} \\
    属格 & \fullpada{pañcānām} & \fullpada{ṣaṇṇām} & \fullpada{saptānām} \\
    依格 & \fullpada{pañcasu} & \fullpada{ṣaṭsu} & \fullpada{saptasu} \\

    & \important{aṣṭa (8)} & \important{nava (9)} & \important{daśa (10)} \\
    主/呼/业 & \fullpada{aṣṭa}/\fullpada{aṣṭau} & \fullpada{nava} & \fullpada{daśa} \\
    具格 & \fullpada{aṣṭabhiḥ}/\fullpada{aṣṭābhiḥ} & \fullpada{navabhiḥ} & \fullpada{daśabhiḥ} \\
    为/从格 & \fullpada{aṣṭabhyaḥ}/\fullpada{aṣṭābhyaḥ} & \fullpada{navabhyaḥ} & \fullpada{daśabhyaḥ} \\
    属格 & \fullpada{aṣṭānām}/\fullpada{aṣṭānām} & \fullpada{navānām} & \fullpada{daśānām} \\
    依格 & \fullpada{aṣṭasu}/\fullpada{aṣṭāsu} & \fullpada{navasu} & \fullpada{daśasu} \\
  \end{NiceTabular}%
  }
\end{frame}


\section{序数词}
\begin{frame}{\insertsection}
  \small
  \tableofcontents[currentsection]
\end{frame}

\begin{frame}{序数词 1-10}
  \small
  \mbox{阳/中性 \nounstem{-a}，阴性 \nounstem{-ā}/\nounstem{-ī}}\\
  \begin{columns}
    \column{0.5\textwidth}
    \begin{itemize}
      \item \nounstem{prathama} 第一
      \item \nounstem{dvitīya} 第二
      \item \nounstem{tṛtīya} 第三
      \item \nounstem{caturtha}/\nounstem{caturīya} 第四
      \item \nounstem{pañcama} 第五
    \end{itemize}
    \column{0.5\textwidth}
    \begin{itemize}
      \item \nounstem{ṣaṣṭha} 第六
      \item \nounstem{saptama} 第七
      \item \nounstem{aṣṭama} 第八
      \item \nounstem{navama} 第九
      \item \nounstem{daśama} 第十
    \end{itemize}
  \end{columns}
\end{frame}


\section{本节作业}

\begin{frame}{\insertsection}
  \begin{itemize}
    \item 第34章阅读练习
  \end{itemize}
\end{frame}

\end{document}
